\documentclass[10pt,a4paper]{amsart}
\usepackage[margin=19mm]{geometry}
\usepackage{amsmath,amssymb,mathtools}
\usepackage[scr=rsfs,cal=euler]{mathalfa}
\usepackage{xfrac}
\usepackage[unicode,hidelinks,bookmarks=false]{hyperref}
\newcommand{\ie}{i.e.\ }
\newcommand{\cf}{cf.\ }
\newcommand{\resp}{respectively\ }
\newcommand{\Subsection}{\subsection}
\providecommand{\nobreakdash}{\nobreak\hbox{-}}
\setlength{\emergencystretch}{2em}
\multlinegap0pt
\usepackage{bm}
\usepackage{bbm}
\usepackage{dsfont}
\usepackage{xspace}
\usepackage{cancel}
\usepackage{mathtools}
\usepackage{amsthm}
\makeatletter
\@ifundefined{subsubsection}{\newtheorem{subsubsection}{Subsubsection}}{}
\@ifundefined{lemma}{\newtheorem{lemma}[subsubsection]{Lemma}}{}
\makeatother
\newcommand\D\Delta
\newcommand\frP{\mathfrak{P}}
\newcommand\ot{\otimes}
\newcommand{\Tr}{\textup{Tr}}
\newcommand{\cohog}[2]{\textup{H}^{#1}({#2})}     % plain group
\newcommand{\wt}[1]{\widetilde{#1}}
\renewcommand\a\alpha
\renewcommand\d\delta
\title[Arithmetic volumes of moduli stacks of Shtukas]{Arithmetic volumes of moduli stacks of Shtukas}
\author{Tony Feng, Zhiwei Yun, Wei Zhang}
\date{Source: arXiv:2601.18557v1 (CC BY 4.0)}
\begin{document}
\maketitle

\noindent\emph{Source and licence.} Arithmetic volumes of moduli stacks of Shtukas. Version arXiv:2601.18557v1 (CC BY 4.0), distributed under the Creative Commons Attribution 4.0 International licence, as recorded in the arXiv licence metadata for this exact version and shown on the abstract page https://arxiv.org/abs/2601.18557v1. Full rights notice: licenses/arxiv-2601.18557-CC-BY-4.0.txt.\par\smallskip

\noindent\textit{Editorial scope.} Counted unit: the Lemma whose source label is \texttt{l:Difg} in the article (source0.tex lines 4441--4447, source label \texttt{l:Difg}), delivered together with the article's own complete original proof of it (source0.tex lines 4448--4541, one \texttt{proof} environment of 94 lines). No mathematics is written, repaired or re-derived: statement and proof are the article's own text line for line. Required context retained uncounted: none. Every definition, display and label of those lines is retained unchanged. Omitted from this package and not redistributed: 1--4440, 4542--5220. Nothing inside a retained excerpt is omitted or altered: every retained body is delivered exactly as the article's lines are written, so no omission receipt of any kind accompanies this package. Citations: the retained excerpts carry no citation command at all, so no bibliography entry of the article is transcribed and no reference is redistributed. Reference adaptation: none was needed and none was made; every reference inside the delivered unit resolves inside it, so no unresolved reference or citation is produced. Printed slips kept literal: no printed word, symbol, hypothesis or number of the counted statement or of its proof was altered. Layout and class: the article's own class is \texttt{amsart} with packages including amsmath, amssymb, amsthm, amsfonts, color, mathrsfs, xcolor, geometry, adjustbox, hyperref, tikz-cd, inputenc, fontenc, tcolorbox; the wrapper is the lane's pinned \texttt{amsart} editorial wrapper at 10pt on A4, which restates the theorem-like environments the retained text needs and adds the article's own macro definitions that the retained text needs (\texttt{\textbackslash D}, \texttt{\textbackslash Tr}, \texttt{\textbackslash a}, \texttt{\textbackslash cohog}, \texttt{\textbackslash d}, \texttt{\textbackslash frP}, \texttt{\textbackslash ot}, \texttt{\textbackslash wt}), with extra wrapper packages (bm, bbm, dsfont, xspace, cancel, mathtools). The delivered numbering is the unit's own. Assets: no retained span contains a figure environment, a graphics-inclusion command, a PDF-page inclusion, a table or a TikZ picture; no image file is copied into this package, the package declares no asset dependency and no image file is redistributed with it. Licence: version arXiv:2601.18557v1 of this article is distributed under the Creative Commons Attribution 4.0 International licence, as recorded in the arXiv licence metadata for this exact version and shown on the abstract page https://arxiv.org/abs/2601.18557v1. A full rights notice is delivered at licenses/arxiv-2601.18557-CC-BY-4.0.txt. Characters: every retained excerpt is pure ASCII, so the delivered engine, XeLaTeX with its default Latin Modern text font, reports no missing character.
\begin{lemma}\label{l:Difg}
For any two homogeneous elements $f,g\in R^W$, and $1\le i\le r$, we have
\begin{equation*}
    D_i(fg)\in (D_1(f),\cdots, D_r(f), D_1(g),\cdots, D_r(g))\subset \wt C^\mu_G
\end{equation*}
where the right side means the ideal generated by the listed elements in $\wt C^\mu_G$.
\end{lemma}

\begin{proof}
    Let $I_{f,g}$ be the ideal generated by $\{D_1(f),\cdots, D_r(f), D_1(g),\cdots, D_r(g)\}$ in $\wt C^\mu_G$. We first prove the statement for $i=1$. To show $D_1(fg)\in I_{f,g}$, we need to show
    \begin{equation}\label{fg1}
        [fg]_1\equiv \sum_{i=1}^r[\Xi_{d+e}]_{1,i}\d_i(fg)\mod I_{f,g}.
    \end{equation}
    Note that
    \begin{equation}\label{d fg}
        \d_i(fg)=\d_i(f)[g]_i+[f]_i\d_i(g)+[c_1(TX)]_i\d_i(f)\d_i(g).
    \end{equation}
    Let $d$ and $e$ be the degrees of $f$ and $g$. Using \eqref{d fg}, \eqref{fg1} is equivalent to
    \begin{equation*}
    [f]_1[g]_1\equiv\sum_{i=1}^r[\Xi_{d+e}]_{1,i}\left(\d_i(f)[g]_i+[f]_i\d_i(g)+[c_1(TX)]_i\d_i(f)\d_i(g)\right)\mod I_{f,g}.
    \end{equation*}
    Now we use the definition of $D_i(f)$ to replace $[f]_i$ above by expressions involving $D_i(f)$ and $\d_{i'}(f)$, it suffices to show
    \begin{eqnarray*}
        &&\sum_{i=1}^r[\Xi_d]_{1,i}\d_i(f)\sum_{i=1}^r[\Xi_e]_{1,i}\d_i(g)\\
        &\equiv& \sum_{i=1}^r[\Xi_{d+e}]_{1,i}\d_i(f)\left(\sum_{i'\ge i}[\Xi_e]_{i,i'}\d_{i'}(g)-\sum_{i'< i}[\Xi_{1-e}]_{i',i}\d_{i'}(g)\right)\\
        &+&\sum_{i=1}^r[\Xi_{d+e}]_{1,i}\left(\sum_{i'\ge i}[\Xi_d]_{i,i'}\d_{i'}(f)-\sum_{i'< i}[\Xi_{1-d}]_{i',i}\d_{i'}(f)\right)\d_i(g)\\
        &+&\sum_{i=1}^r[\Xi_{d+e}]_{1,i}[c_1(TX)]_i\d_i(f)\d_i(g)\mod I_{f,g}.
    \end{eqnarray*}
    Comparing coefficients of $\d_i(f)\d_{i'}(g)$ on both sides, it reduces to showing
    \begin{eqnarray}\label{Xi prod}
        [\Xi_d]_{1,i}[\Xi_e]_{1,i'}=\begin{cases}
            [\Xi_{d+e}]_{1,i}[\Xi_{e}]_{i,i'}-[\Xi_{d+e}]_{1,i'}[\Xi_{1-d}]_{i,i'} & i<i'\\
            -[\Xi_{d+e}]_{1,i}[\Xi_{1-e}]_{i',i}+[\Xi_{d+e}]_{1,i'}[\Xi_d]_{i',i}& i>i'\\
            [\Xi_{d+e}]_{1,i}\left([\Xi_e]_{i,i}+[\Xi_d]_{i,i}+[c_1(TX)]_i\right) & i=i'.
        \end{cases}
    \end{eqnarray}
    Let us prove the above identity in the case $i<i'$. Write $x=q^{d-1}$ and $y=q^{e-1}$, the identity can be written as
    \begin{eqnarray}\label{prod D phi}
        &&\left(1\ot \frac{1}{\phi x-1}\right)[\D]_{1,i}\left(1\ot \frac{1}{\phi y-1}\right)[\D]_{1,i'}\\
\notag        &=& \left(1\ot \frac{1}{q\phi xy-1}\right)[\D]_{1,i}\left(1\ot \frac{1}{\phi y-1}\right)[\D]_{i,i'}\\
\notag        &+&\left(1\ot \frac{1}{q\phi xy-1}\right)[\D]_{1,i'}\left(\frac{\phi x}{\phi x-1}\ot1\right)[\D]_{i,i'}.
    \end{eqnarray}
    Expand both sides in geometric series in $x$ and $y$, the coefficient of $x^ay^b$ on the left side is 
    \begin{equation}\label{left xy}
        (1\ot \phi^a)[\D]_{1,i}(1\ot \phi^b)[\D]_{1,i'}=(1\ot \phi^a\ot \phi^b)[\D]_{1,i,i'}.
    \end{equation}
    Here we use that $[\D]_{1,i}[\D]_{1,i'}$ is the class of the small diagonal in $X\times X\times X$ (the $1,i,i'$-factors), which we denote by $[\D]_{1,i,i'}$. The coefficient of $x^ay^b$ on the right side of \eqref{prod D phi} is
    \begin{equation}\label{right xy}
        \begin{cases}
            (1\ot q^b\phi^b)[\D]_{1,i'}(\phi^{a-b}\ot1)[\D]_{i,i'}& a>b\\
            (1\ot q^a\phi^a)[\D]_{1,i}(1\ot \phi^{b-a})[\D]_{i,i'}  & a\le b.
        \end{cases}
    \end{equation}
    When $a>b$, we have
    \begin{equation*}
        (1\ot q^b\phi^b)[\D]_{1,i'}(\phi^{a-b}\ot1)[\D]_{i,i'}=q^b(1\ot \phi^b)[\D]_{1,i'}(\phi^{a-b}\ot1)[\D]_{i,i'}=(1\ot \phi^b)[\D]_{1,i'}(\phi^{a-b}\ot1)(\phi\ot \phi)^b[\D]_{i,i'}.
    \end{equation*}
    Here we use that $(\phi\ot\phi)[\D]_{i,i'}=\phi[\D]_{i,i'}=q[\D]_{i,i'}$, because $[\D]$ is a divisor class. Therefore
    \begin{equation*}
        (1\ot q^b\phi^b)[\D]_{1,i'}(\phi^{a-b}\ot1)[\D]_{i,i'}=(1\ot \phi^b)[\D]_{1,i'}(\phi^a\ot \phi^{b})[\D]_{i,i'}=(1\ot \phi^a\ot \phi^b)[\D]_{1,i,i'},
    \end{equation*}
    which is equal to \eqref{left xy}.
    Here we use that $\phi^a$ is a ring endomorphism of $\cohog{*}{X}$. This shows that \eqref{right xy} and \eqref{left xy} are equal when $a>b$. The proof of their equality when $a\le b$ is similar. Thus we have verified the $i<i'$ case of \eqref{Xi prod}.

    The case $i>i'$ of \eqref{Xi prod} is proved similarly. 

    Finally let us prove the $i=i'$ case of \eqref{Xi prod}. When $i=1=i'$ both sides are zero. We therefore assume $i=i'>1$. Then the equality \eqref{Xi prod} becomes
    \begin{equation}\label{ii coeff}
        \left(1\ot \frac{1}{\phi x-1}\right)[\D]_{1,i}\left(1\ot \frac{1}{\phi y-1}\right)[\D]_{1,i}=\left(1\ot \frac{1}{q\phi xy-1}\right)[\D]_{1,i}\cdot \Tr\left(\frac{1}{\phi x-1}+\frac{1}{\phi y-1}+1\Big|\cohog{*}{X}\right)[\xi]_i,
    \end{equation}
    Since $\phi$ and $q\phi^{-1}$ are adjoint under the Poincar\'e duality pairing on $\cohog{*}{X}$, we have
    \begin{equation*}
    \left(1\ot \frac{1}{\phi y-1}\right)[\D]_{1,i}=\left(\frac{1}{q\phi^{-1}y-1}\ot 1\right)[\D]_{1,i}.
    \end{equation*}
    Therefore the left side of \eqref{ii coeff} is
    \begin{equation*}
    \left(1\ot \frac{1}{\phi x-1}\right)[\D]_{1,i}\left(\frac{1}{q\phi^{-1}y-1}\ot 1\right)[\D]_{1,i}=\Tr\left(\frac{1}{q\phi^{-1}y-1}\cdot\frac{1}{\phi x-1}\Big|\cohog{*}{X}\right)[\xi\ot\xi]_{1,i}.
    \end{equation*}
    We also have
    \begin{equation*}
    \left(1\ot \frac{1}{q\phi xy-1}\right)[\D]_{1,i}\cdot [\xi]_i=\frac{1}{qxy-1}[\xi\ot\xi]_{1,i}.
    \end{equation*}
    Therefore \eqref{ii coeff} is equivalent to the equality
    \begin{equation}\label{trace frac phi}
    \Tr\left(\frac{1}{q\phi^{-1}y-1}\cdot\frac{1}{\phi x-1}\right)=\frac{1}{qxy-1}\Tr\left(\frac{1}{\phi x-1}+\frac{1}{\phi y-1}+1\right).
    \end{equation}
    where $\Tr$ means alternating trace on $\cohog{*}{X}$. Again we may change $\frac{1}{\phi y-1}$ on the right side above by its adjoint $\frac{1}{q\phi^{-1}y-1}$, and \eqref{trace frac phi} follows from the equality of endomorphisms of $\cohog{*}{X}$ before taking trace
    \begin{equation*}
    \frac{1}{q\phi^{-1}y-1}\cdot\frac{1}{\phi x-1}=\frac{1}{qxy-1}\left(\frac{1}{\phi x-1}+\frac{1}{q\phi^{-1} y-1}+1\right).
    \end{equation*}
    This finishes the proof of the lemma for $i=1$.

    Now consider the case of general $i$. Let $\mu'=(\mu_2,\mu_3,\cdots, \mu_r,\mu_1)$ (so that $\mu'_i=\mu_{i+1}$, with the subscripts understood mod $r$). We define a partial Frobenius
    \begin{equation}\label{partial Frob on C}
        \frP^*: \wt C^{\mu'}_G\to \wt C^{\mu}_G
    \end{equation}
    that maps $\a_1\ot\a_2\ot\cdots\ot \a_r$ to $\phi(\a_r)\ot\a_1\ot\cdots\ot \a_{r-1}$, where $\a_i\in \cohog{*}{X\times G/P_{\mu'_{i}}}$. It is easy to check that $D_2(f)=\frP^*(D_1(f))\in \wt C^\mu_G$, where $D_1(f)$ is viewed as an element of $\wt C^{\mu'}_G$. We have proved that $D_1(fg)\in I'_{f,g}:=(D_1(f),\cdots, D_r(f), D_1(g),\cdots, D_r(g))\subset \wt C^{\mu'}_G$. Applying $\frP^*$, we see that
    \begin{equation*}
        D_2(fg)=\frP^*D_1(fg)\in \frP^* I'_{f,g}=I_{f,g}\subset \wt C^\mu_G.
    \end{equation*}
    Repeating this argument we see that $D_i(fg)\in I_{f,g}$ for all $i=2,\cdots, r$. This finishes the proof of the lemma.
\end{proof}

\end{document}
